\chapter{Composing a Driver, End to End}
\label{ch:composing}

\chapterorient{Part~V disassembled the simulator into named machines and showed,
in a drawn trace (\cref{ch:fulltrace}), that one analysis is one composition of
them. Part~VI takes that picture seriously: if every analysis is a composition of
the same library functions, then the compositions are \emph{code}, and we can
write them down. This chapter does exactly that. We open the \code{drivers}
library---the topmost shelf of the synthesized library (\cref{ch:fivelibs})---and
read its six driver definitions, each a handful of lines over the four calculus
libraries beneath it. We render one driver, electrostatics, in full and walk it
line by line, naming the library every call comes from. Then we descend into the
single most important step in the book, the conjugate-gradient \code{krylov\_step},
and watch the abstract calculus of Part~III produce real, runnable-looking
algorithm code. The payoff is the chapter's thesis: every driver is a few lines of
\emph{map}, \emph{fold}, or \emph{single call} over \texttt{assemble}$\to$%
\texttt{reduce}, and the calculus we built really does synthesize working
algorithms.}

\section{The \texttt{drivers} library: six compositions over four}

We have, by now, the whole library. \Cref{ch:synthlib} laid out its
architecture---five modules---and \cref{ch:fivelibs} toured the four that hold the
\emph{vocabulary}: \code{types} (the records every solve threads), \code{iteration}
(the loop combinators of \cref{ch:fold} and the step kernels), \code{data-algebra}
(the assembly fold \code{fe\_assemble} and the reduction verbs \code{gram\_reduce},
\code{sparameter\_reduce}, and their siblings), and \code{coordination} (the
outer-driver shapes \code{solve\_family}, \code{frequency\_sweep},
\code{fold\_solve}, and the opaque \code{eigsolve} cap). This chapter is about the
\emph{fifth} module, the one that sits on top of all of them:

\begin{quote}\itshape
The \code{drivers} library holds the six entry-point analyses---electrostatic,
magnetostatic, driven, transient, eigenmode, boundary-mode---each written as a short
composition of the four libraries below it. It composes the vocabulary; it does not
add to it.
\end{quote}

\noindent That is the whole nature of the module, and it is worth stating sharply
before we read any code. A driver definition introduces no new operator, no new
combinator, no new record. Every name in its body has already been built and named
in an earlier chapter; the driver's only job is to \emph{wire those names together}
in the right shape. Reading the \code{drivers} library is therefore reading the
book's table of contents executed: each line points down at a machine you already
own (\cref{fig:lib-stack}).

\begin{figure}[t]
\centering
\begin{tikzpicture}[font=\footnotesize, >=Stealth]
  % top: the drivers module
  \node[stage, fill=simBgViolet, draw=simViolet, minimum width=92mm,
        minimum height=11mm, align=center] (drv) at (0,3.2)
    {\textbf{\code{drivers}}\ \scriptsize --- the 6 entry-point compositions
     (this chapter)};
  % the four calculus libraries below
  \node[stage, fill=simBgGold, draw=simGold, minimum width=20mm,
        minimum height=15mm, align=center] (da) at (-3.45,1.1)
    {\code{data-}\\\code{algebra}\\[1pt]\scriptsize assemble\\+ reduce};
  \node[stage, fill=simBgRust, draw=simRust, minimum width=20mm,
        minimum height=15mm, align=center] (co) at (-1.15,1.1)
    {\code{coord-}\\\code{ination}\\[1pt]\scriptsize the solve\\shapes};
  \node[stage, fill=simBgBlue, draw=simBlue, minimum width=20mm,
        minimum height=15mm, align=center] (it) at (1.15,1.1)
    {\code{iter-}\\\code{ation}\\[1pt]\scriptsize loop +\\step};
  \node[stage, fill=simBgGray, draw=simGray, minimum width=20mm,
        minimum height=15mm, align=center] (ty) at (3.45,1.1)
    {\code{types}\\[1pt]\scriptsize the\\records};
  % composition edges
  \foreach \n in {da,co,it,ty}{ \draw[depedge] (drv.south) -- (\n.north); }
  \node[font=\scriptsize\itshape, text=simGray] at (0,0.05)
    {(each library composes the ones to its right)};
\end{tikzpicture}
\caption[The library stack]{The synthesized library as a stack. The
\code{drivers} module (violet, this chapter) sits on top and \emph{composes} the
four calculus libraries of \cref{ch:fivelibs}: \code{data-algebra} (the assembly
fold and the reduction verbs), \code{coordination} (the four solve shapes plus the
opaque eigensolver cap), \code{iteration} (the loop combinators and step kernels),
and \code{types} (the records). A driver definition adds no new vocabulary; every
name in its body is built below it. Reading a driver is reading those four
libraries wired into the shape one analysis needs.}
\label{fig:lib-stack}
\end{figure}

\subsection{The unifying algebra, in one line each}

Here are all six drivers, each collapsed to the composition that defines it. Read
each right to left, the way function composition reads---\emph{assemble} on the
right, \emph{reduce} on the left, the solve shape in the middle:

\begin{lstlisting}[style=l4style]
electrostatic = gram_reduce(w=1)          . solve_family    . fe_assemble
magnetostatic = gram_reduce(w=1/(Ii*Ij))  . solve_family    . fe_assemble
driven        = sparameter_reduce         . frequency_sweep . fe_assemble x3
transient     = fold_solve                                  . fe_assemble x3
eigenmode     = map readout . eigsolve    . eig_pencil      . fe_assemble x3
boundary_mode = map readout . eigsolve    . eig_pencil      . fe_assemble . extract_2d
\end{lstlisting}

\noindent Six analyses, six one-liners, and one shape. Every driver is
\texttt{assemble} on the right, a \emph{solve shape} in the middle, and a
\texttt{reduce} on the left. What differs from row to row is small and named: how
many operators get assembled, which of the four \emph{coordination shapes} runs the
solve, and which reduction verb caps it. \Cref{ch:situating} taught us to read those
middle terms as a small grid of \emph{solve corners}; the \code{drivers} library is
that grid made executable. \Cref{fig:six-trees} draws the six as composition trees,
grouped by the corner each occupies.

\begin{figure}[p]
\centering
\begin{tikzpicture}[font=\scriptsize, >=Stealth,
    every node/.style={align=center}]
% ===================== GROUP 1: the maps =====================
\node[cornertag, anchor=west, font=\footnotesize\bfseries] at (-0.3,7.7)
  {\textsc{the maps} --- a family of independent solves (\cref{ch:situating}, corners 1--2)};

% electrostatic
\node[stage, fill=simBgGold, draw=simGold, text width=20mm] (e0) at (1.4,6.4)
  {\code{electrostatic}};
\node[opbox, text width=15mm] (e1) at (-0.4,5.2) {\code{fe\_assemble}\\\scriptsize K once};
\node[stage, fill=simBgBlue, draw=simBlue, text width=15mm] (e2) at (1.4,5.2)
  {\code{solve\_family}\\\scriptsize fixed-op map};
\node[opbox, text width=15mm] (e3) at (3.2,5.2) {\code{gram\_reduce}\\\scriptsize $w{=}1$};
\draw[depedge] (e0)--(e1); \draw[depedge] (e0)--(e2); \draw[depedge] (e0)--(e3);

% magnetostatic
\node[stage, fill=simBgGold, draw=simGold, text width=20mm] (m0) at (6.6,6.4)
  {\code{magnetostatic}};
\node[opbox, text width=15mm] (m1) at (4.8,5.2) {\code{fe\_assemble}\\\scriptsize K once};
\node[stage, fill=simBgBlue, draw=simBlue, text width=15mm] (m2) at (6.6,5.2)
  {\code{solve\_family}\\\scriptsize fixed-op map};
\node[opbox, text width=16mm] (m3) at (8.5,5.2) {\code{gram\_reduce}\\\scriptsize $w{=}\frac{1}{I_iI_j}$};
\draw[depedge] (m0)--(m1); \draw[depedge] (m0)--(m2); \draw[depedge] (m0)--(m3);

% driven
\node[stage, fill=simBgGold, draw=simGold, text width=18mm] (d0) at (11.4,6.4)
  {\code{driven}};
\node[opbox, text width=15mm] (d1) at (9.7,5.2) {\code{fe\_assemble}\\\scriptsize $\times3$ basis};
\node[stage, fill=simBgRust, draw=simRust, text width=16mm] (d2) at (11.4,5.2)
  {\code{frequency\_}\\\code{sweep}\\\scriptsize op-varying map};
\node[opbox, text width=16mm] (d3) at (13.3,5.2) {\code{sparameter\_}\\\code{reduce}};
\draw[depedge] (d0)--(d1); \draw[depedge] (d0)--(d2); \draw[depedge] (d0)--(d3);

% ===================== GROUP 2: the fold =====================
\node[cornertag, anchor=west, font=\footnotesize\bfseries] at (-0.3,3.9)
  {\textsc{the fold} --- a state-threaded chain (\cref{ch:situating}, corner 3)};
\node[stage, fill=simBgGold, draw=simGold, text width=18mm] (t0) at (2.0,2.7)
  {\code{transient}};
\node[opbox, text width=16mm] (t1) at (0.0,1.5) {\code{fe\_assemble}\\\scriptsize $\times3$ K,C,M};
\node[stage, fill=simBgGreen, draw=simGreen, text width=18mm] (t2) at (2.6,1.5)
  {\code{fold\_solve}\\\scriptsize state-threaded\\time march};
\draw[depedge] (t0)--(t1); \draw[depedge] (t0)--(t2);

% ===================== GROUP 3: the black box =====================
\node[cornertag, anchor=west, font=\footnotesize\bfseries] at (6.2,3.9)
  {\textsc{the black box} --- one opaque call (\cref{ch:situating}, corner 4)};

% eigenmode
\node[stage, fill=simBgGold, draw=simGold, text width=16mm] (g0) at (8.0,2.7)
  {\code{eigenmode}};
\node[opbox, text width=14mm] (g1) at (6.6,1.5) {\code{fe\_assemble}\\\scriptsize $\times3$ pencil};
\node[extern, text width=12mm] (g2) at (8.1,1.5) {\code{eigsolve}\\\scriptsize 1 call};
\node[opbox, text width=12mm] (g3) at (9.5,1.5) {\code{map}\\\code{readout}};
\draw[depedge] (g0)--(g1); \draw[depedge] (g0)--(g2); \draw[depedge] (g0)--(g3);

% boundary mode
\node[stage, fill=simBgGold, draw=simGold, text width=18mm] (b0) at (12.4,2.7)
  {\code{boundary\_mode}};
\node[opbox, text width=14mm] (b1) at (10.7,1.5) {\code{extract\_2d}\\\scriptsize submesh};
\node[opbox, text width=13mm] (b2) at (12.1,1.5) {\code{fe\_assemble}\\\scriptsize $+$ pencil};
\node[extern, text width=11mm] (b3) at (13.5,1.5) {\code{eigsolve}\\\scriptsize 1 call};
\draw[depedge] (b0)--(b1); \draw[depedge] (b0)--(b2); \draw[depedge] (b0)--(b3);

\end{tikzpicture}
\caption[The six drivers as composition trees]{The hero figure: all six drivers as
composition trees, grouped by the solve corner each occupies (\cref{ch:situating}).
\textsc{The maps} (top): electrostatic and magnetostatic run \code{solve\_family},
a fixed-operator map (corner~1); the driven sweep runs \code{frequency\_sweep}, an
operator-varying map (corner~2). \textsc{The fold} (lower left): transient threads
its state through \code{fold\_solve}, a sequential chain (corner~3). \textsc{The
black box} (lower right): eigenmode and boundary-mode hand the assembled pencil to a
single opaque \code{eigsolve} call (corner~4), boundary-mode prefacing it with a
2-D submesh extraction. Every tree has the same skeleton---an assemble on one side,
a reduce/readout on the other, a solve shape between---and differs only in which
shape and which reduction it picks. This grid \emph{is} the \code{drivers} library.}
\label{fig:six-trees}
\end{figure}

\subsection{The three coordination shapes are the deep unifier}

Look again at the middle column of the six one-liners and a smaller pattern
emerges---smaller even than the four corners of \cref{ch:situating}. Underneath the
corners there are really only \emph{three coordination shapes}, the three ways a
driver can run its solves (\cref{ch:combinators}):

\begin{description}[style=nextline, leftmargin=0pt, font=\normalfont\itshape]
  \item[map] a family of \emph{independent} solves, run in any order---no carry.
  Electrostatics and magnetostatics map over a fixed operator (\code{solve\_family});
  the driven sweep maps over a per-frequency operator (\code{frequency\_sweep}). The
  shape is the same \texttt{map}; the corners differ only in whether the operator is
  reused or rebuilt.
  \item[fold] a \emph{sequential} chain, each step consuming the state the last
  produced. Transient time-stepping is the canonical fold (\code{fold\_solve}); so is
  the adaptive mesh loop that the lifecycle wraps around all of it.
  \item[single call] no loop \emph{we} write at all---the assembled pencil goes to an
  opaque library routine (\code{eigsolve}) and the modes come back. Eigenmode and
  boundary-mode are the two black-box drivers.
\end{description}

\noindent Map, fold, single call. That is the entire vocabulary of \emph{how a
driver coordinates its solves}, and every analysis Palace runs is one of the three
wrapped around \texttt{assemble} and capped by a \texttt{reduce}. \Cref{fig:three-shapes}
sets them side by side; the rest of the chapter is, in effect, three readings of
this one figure---one map (electrostatics, in full), one fold (the CG step that
drives every solve inside any of them), and a tour of the single call.

\begin{figure}[t]
\centering
\begin{tikzpicture}[font=\footnotesize, >=Stealth]
  % ---- MAP ----
  \begin{scope}[shift={(0,0)}]
    \node[font=\small\bfseries, text=simBlue] at (1.3,2.2) {\code{map}};
    \node[opbox, minimum width=10mm] (mk) at (1.3,1.55) {op (fixed)};
    \foreach \j/\y in {1/0.7,2/0.0,3/-0.7}{
      \node[data, minimum width=7mm] (mb\j) at (0.3,\y) {$\vb_{\j}$};
      \node[product, minimum width=8mm] (mv\j) at (2.3,\y) {$\vx_{\j}$};
      \draw[flow, shorten >=1pt] (mb\j) -- (mv\j);
    }
    \node[font=\scriptsize\itshape, text=simBlue, align=center] at (1.3,-1.4)
      {independent;\\reorderable, parallel};
  \end{scope}
  \draw[simGray!45, thick] (3.6,-1.7) -- (3.6,2.4);
  % ---- FOLD ----
  \begin{scope}[shift={(4.6,0)}]
    \node[font=\small\bfseries, text=simGreen] at (1.6,2.2) {\code{fold}};
    \node[data, minimum width=8mm] (s0) at (-0.1,0.7) {$\vs_0$};
    \node[opbox, minimum width=7mm] (st1) at (0.9,0.7) {step};
    \node[data, minimum width=8mm] (s1) at (1.85,0.7) {$\vs_1$};
    \node[opbox, minimum width=7mm] (st2) at (2.8,0.7) {step};
    \node[data, minimum width=8mm] (s2) at (3.75,0.7) {$\vs_2$};
    \draw[flow] (s0)--(st1); \draw[flow] (st1)--(s1);
    \draw[flow] (s1)--(st2); \draw[flow] (st2)--(s2);
    \node[font=\scriptsize\itshape, text=simGreen, align=center] at (1.8,-1.4)
      {sequential;\\each step needs the last};
  \end{scope}
  \draw[simGray!45, thick] (9.0,-1.7) -- (9.0,2.4);
  % ---- SINGLE CALL ----
  \begin{scope}[shift={(9.9,0)}]
    \node[font=\small\bfseries, text=simViolet] at (1.5,2.2) {single call};
    \node[data, minimum width=12mm] (pen) at (0.2,0.7) {pencil};
    \node[extern, minimum width=15mm] (eig) at (2.2,0.7) {\code{eigsolve}};
    \node[product, minimum width=12mm] (modes) at (2.2,-0.4) {modes};
    \draw[flow] (pen) -- (eig);
    \draw[flow] (eig) -- (modes);
    \node[font=\scriptsize\itshape, text=simViolet, align=center] at (1.5,-1.4)
      {opaque; no loop\\we write};
  \end{scope}
\end{tikzpicture}
\caption[The three coordination shapes]{The three ways a driver coordinates its
solves---the deep unifier beneath the six drivers. A \code{map} (blue) runs a family
of independent solves over an operator; nothing is carried, so the solves are
reorderable and parallel (electrostatic, magnetostatic, driven). A \code{fold}
(green) threads a state through a sequential chain, each step consuming the last
(transient, and the adaptive mesh loop). A \code{single call} (violet) hands the
assembled pencil to an opaque library routine and reads the modes back (eigenmode,
boundary-mode). Every driver in the library is one of these three, wrapped around
\code{fe\_assemble} and capped by a reduction.}
\label{fig:three-shapes}
\end{figure}

\begin{keyidea}
Every driver is a few lines: a \texttt{map}, a \texttt{fold}, or a single opaque
call, wrapped around \texttt{assemble} and capped by \texttt{reduce}. The six
analyses differ only in \emph{which} of the three coordination shapes runs the solve
and \emph{which} reduction verb caps it---and the calculus of Part~III produces real
algorithm code that fills those shapes in. The \code{drivers} library is the grid of
\cref{ch:situating} written out as composition; reading it is reading the book's
machinery wired together.
\end{keyidea}

\subsection{Per-driver config records are projections, not new shapes}
\label{sec:config-projections}

Before we render a driver in full, one piece of plumbing deserves naming, because it
recurs in every signature. Each driver's body opens by reading fields off a
configuration value typed \code{ElectrostaticConfig}, \code{DrivenConfig}, and so on.
These look like six different record types. They are not. There is exactly
\emph{one} configuration record---the \code{IoData} parsed from the user's file
(\cref{ch:records}, \cref{ch:fulltrace})---and the per-driver config types are
\emph{projection-views} of it: type aliases that name which sub-records a given
driver reads.

\begin{lstlisting}[style=l4style]
-- the per-driver config records are VIEWS of the one IoData, not new shapes
type ElectrostaticConfig = IoData   -- reads: model, domains.epsilon, terminals, linear solver
type DrivenConfig        = IoData   -- reads: model, domains.{nu,sigma,eps}, ports, omega-sweep
type EigenmodeConfig     = IoData   -- reads: model, domains.{nu,sigma,eps}, n-modes + target
-- ... and likewise for magnetostatic, transient, boundary_mode
\end{lstlisting}

\noindent Writing \code{ElectrostaticConfig} rather than the bare \code{IoData} in a
driver's type is documentation, not a new data shape: it says ``this driver consults
\emph{this} slice of the one config.'' The accessors a driver uses---%
\lstinline[style=l4style]|permittivity cfg|, \lstinline[style=l4style]|terminal_sources cfg|,
\lstinline[style=l4style]|ports cfg|---are the projection's utility API, trivial
field selectors into \code{IoData}. Keeping these as views and not as separate
records is the record-definition discipline of \cref{ch:records} applied at the top
of the stack: the data shape is defined once, in \code{types}, and the drivers
\emph{view} it rather than redefining it.

\begin{notationbox}
The shared \code{IoData} is why the lifecycle (\cref{ch:fulltrace}) can parse one
file and dispatch to any of the six drivers: the parse produces a single
\code{IoData}, and \emph{dispatch} reads one field of it---%
\lstinline[style=l4style]|problemType cfg|---to choose which driver to run. The
chosen driver then projects the slice it needs. One parse, one record, six views, a
one-field switch: that is the entire shared front end, and it is why
\S\ref{sec:lifecycle} can write the spine root as a dispatch over the six driver
defs.
\end{notationbox}

\section{One driver, rendered and walked: \texttt{electrostatic}}
\label{sec:es-rendered}

We now render a complete driver as synthesized code and read every line of it. We
choose \code{electrostatic}---the cleanest of the six, the anchor of all of Part~V
(\cref{ch:electrostatics})---because its shape is the bare skeleton: assemble one
operator, map a family of solves over it, Gram-reduce the family to the answer. Once
we have walked this one, the other five are short variations, and we contrast the
richest of them, \code{driven}, in \S\ref{sec:driven-contrast}.

\subsection{The whole driver}

Here is the synthesized \code{electrostatic} definition in full---the rendered form
of the composition $\code{gram\_reduce}\circ\code{solve\_family}\circ
\code{fe\_assemble}$ that \cref{ch:electrostatics} anchored Part~V on:

\begin{lstlisting}[style=l4style]
-- # Arguments
--   cfg : ElectrostaticConfig   -- the IoData electrostatic view: H1 space, epsilon,
--                                  terminal sources, linear solver config
-- # Returns
--   CapacitanceMatrix           -- the n x n Maxwell capacitance matrix (the product)
electrostatic :: ElectrostaticConfig -> CapacitanceMatrix
electrostatic cfg =
  let space = h1_space cfg                                  -- (0) the H1 FE space
      k     = fe_assemble space [ diffusion (permittivity cfg) ]   -- (1) assemble K ONCE
      rhss  = [ excitation cfg idx | idx <- terminal_sources cfg ] -- (2) per-terminal RHS family
      vs    = solve_family k rhss                           -- (3) fixed-operator solve map
  in  gram_reduce k vs (\i j -> 1)                          -- (4) C_ij = V_j^T K V_i, w = 1
\end{lstlisting}

\noindent Five lines of body. That is a complete electromagnetic analysis---the same
one that, on a real mesh, fills a capacitance matrix for a chip interconnect. The
brevity is the point: the driver is \emph{nothing but} the composition. All the
weight lives in the four library calls, each built in an earlier chapter. We walk
them in order, naming the library each comes from (\cref{fig:es-dataflow}).

\subsection{Line by line, naming the library}
\label{sec:es-linewalk}

\paragraph{(0) \code{space = h1\_space cfg}---from \code{types}/\code{data-algebra}.}
The first line builds the finite-element space the analysis lives in: the scalar
nodal $\Hone$ space (\cref{ch:electrostatics}), read off the mesh and polynomial
order in \code{cfg}. It is construction-stratum data (\cref{ch:strata})---fixed the
instant the mesh is built, read-only thereafter. Everything downstream is expressed
against this space.

\paragraph{(1) \code{k = fe\_assemble space [\,diffusion (permittivity cfg)\,]}---from
\code{data-algebra}.} The second line assembles the operator, and it is the
load-bearing line of the whole fixed-operator corner. \code{fe\_assemble} is the
assembly fold of \cref{ch:assembly}: it walks a list of weak-form terms and sums one
contribution per term into a sparse operator. For electrostatics the list has
\emph{exactly one} term---the $\varepsilon$-weighted gradient,
\code{diffusion (permittivity cfg)}, the discrete form of
$\int_\dom\varepsilon\Grad u\cdot\Grad v\dV$ (\cref{ch:electrostatics},
eq.~\eqref{eq:es-bilinear-ch39})---so the fold runs once and returns the stiffness
operator $\Kop$. The word that matters is the one the type makes invisible:
\emph{once}. This line sits \emph{outside} the solve map below it, so $\Kop$ is built
a single time and reused for every terminal. Hoisting it here, structurally, is the
fixed-operator corner (\cref{ch:electrostatics}, and the pitfall it warned of).

\paragraph{(2) \code{rhss = [\,excitation cfg idx | idx <- terminal\_sources cfg\,]}---from
the config view.} The third line builds the family of right-hand sides, one per
terminal. The list comprehension ranges over \code{terminal\_sources cfg}---the
conductors named in the config---and for each builds the vector
\code{excitation cfg idx} that encodes ``conductor \code{idx} at one volt, the rest
grounded'' (the eliminated boundary data, \cref{ch:electrostatics}). This is the
\emph{method of unit excitations} (\cref{ch:electrostatics}) rendered as a
comprehension: $n$ conductors, $n$ right-hand sides, no solving yet.

\paragraph{(3) \code{vs = solve\_family k rhss}---from \code{coordination}.} The
fourth line runs the solves, and it is the coordination shape of the driver. The
\code{solve\_family} combinator (\cref{ch:combinators}) is the \emph{fixed-operator
map}: it captures the operator \code{k} once and maps a solver over the RHS family,
\lstinline[style=l4style]|map (ksp_solve k) rhss|. Because \code{k} is a free
variable of the mapped function, the \texttt{map} closes over it once and reuses it
for every element---the hoist made structural by the type, so no implementation can
rebuild the operator inside the loop. Each element solve \code{ksp\_solve k b\_j} is
one preconditioned conjugate-gradient solve of the SPD system $\Kop V_j=\vb_j$, and
\code{ksp\_solve} itself bottoms out in the \code{krylov\_step} of
\S\ref{sec:cg-flagship}---the step we render in full next. The result \code{vs} is
the solution family $[V_1,\dots,V_n]$.

\paragraph{(4) \code{gram\_reduce k vs (\textbackslash i j -> 1)}---from
\code{data-algebra}.} The final line reduces the family to the answer.
\code{gram\_reduce} (\cref{ch:reductions}) is the symmetric Gram verb: it pairs the
family with itself through the operator \code{k}, computing
$C_{ij}=V_j^{\!\top}\Kop V_i$ over the upper triangle and mirroring the lower (since
$\Kop$ is symmetric). The weight closure \lstinline[style=l4style]|\i j -> 1| is the
unit-voltage specialization that makes this reduction \emph{capacitance}
specifically; magnetostatics will pass a different weight to the same verb
(\S\ref{sec:driven-contrast}). Notice that \code{k} appears \emph{twice}---in the
solve (line 3) and again in the reduce (line 4): the same operator that drives the
solves also weights the Gram. That double use is the back-edge of the anchor diagram
(\cref{ch:electrostatics}, \cref{fig:es-pipeline}), now visible in the code as a
shared \code{let}-binding.

\begin{figure}[t]
\centering
\begin{tikzpicture}[node distance=7mm and 11mm, font=\footnotesize, >=Stealth]
  \node[data, align=center, text width=17mm] (cfg)
    {\code{cfg}\\[1pt]\scriptsize \code{IoData} view};
  \node[stage, right=of cfg, align=center, text width=20mm, fill=simBgGold, draw=simGold]
    (asm) {\code{fe\_assemble}\\[1pt]\scriptsize (1) $\Kop$ once};
  \node[stage, right=of asm, align=center, text width=22mm, fill=simBgBlue, draw=simBlue]
    (solve) {\code{solve\_family}\\[1pt]\scriptsize (3) fixed-op map};
  \node[stage, right=of solve, align=center, text width=20mm, fill=simBgGold, draw=simGold]
    (red) {\code{gram\_reduce}\\[1pt]\scriptsize (4) $w{=}1$};
  \node[product, right=of red, align=center, text width=15mm] (prod)
    {$\mat{C}$\\[1pt]\scriptsize capacitance};
  \draw[flow] (cfg) -- (asm);
  \draw[flow] (asm) -- node[above,font=\scriptsize]{$\Kop$} (solve);
  % the RHS family entering the solve from below
  \node[data, below=8mm of solve, align=center, text width=20mm] (rhs)
    {\code{rhss} $=[\vb_j]$\\[1pt]\scriptsize (2) per terminal};
  \draw[flow] (rhs.north) -- (solve.south);
  \draw[flow] (solve) -- node[above,font=\scriptsize]{$\{V_j\}$} (red);
  \draw[flow] (red) -- (prod);
  % the shared K back-edge
  \draw[depedge, shorten >=1pt] (asm.south) ++(0,-1mm)
    .. controls +(0,-12mm) and +(0,-12mm) .. (red.south)
    node[midway, below, font=\scriptsize, text=simBlue]
      {the \emph{same} $\Kop$ weights the Gram (\code{k} used twice)};
  % library tags
  \node[cornertag, above=0.5mm of asm]   {\code{data-algebra}};
  \node[cornertag, above=0.5mm of solve] {\code{coordination}};
  \node[cornertag, above=0.5mm of red]   {\code{data-algebra}};
\end{tikzpicture}
\caption[The \texttt{electrostatic} driver's data flow]{The data flow of the
rendered \code{electrostatic} driver (\S\ref{sec:es-rendered}), each stage tagged by
the library its call comes from. The config view \code{cfg} feeds
\code{fe\_assemble} (\code{data-algebra}), which builds the stiffness $\Kop$
\emph{once}; the per-terminal RHS family \code{rhss} enters from below;
\code{solve\_family} (\code{coordination}) maps the fixed-operator solve over the
family, yielding $\{V_j\}$; and \code{gram\_reduce} (\code{data-algebra}) pairs that
family against $\Kop$ into the capacitance matrix $\mat{C}$. The blue back-edge is
the same operator $\Kop$ consumed twice---by the solve and again by the reduce---which
is the \code{k} that appears in both line~(3) and line~(4) of the code.}
\label{fig:es-dataflow}
\end{figure}

\subsection{What the rendering shows that the diagram could not}

The trace of \cref{ch:fulltrace} \emph{drew} this pipeline; the rendering
\emph{types} it. Three things become concrete only in the code. First, the operators
flow as ordinary \code{let}-bound values---$\Kop$ is a value named \code{k}, passed
to two later lines---so ``the same operator weights the Gram'' is not a diagram
annotation but a variable used twice. Second, the per-terminal family is a list
comprehension, so ``$n$ solves'' is a structural consequence of
\code{terminal\_sources cfg} having $n$ elements, not a loop bound written by hand.
Third, the whole driver has the type \code{ElectrostaticConfig -> CapacitanceMatrix}:
config in, product out, no monad, no mutation, no mention of the iteration machinery
that the solve hides. The driver is pure all the way across; the only place the
state monad of \cref{ch:monad} appears is buried four levels down, inside
\code{ksp\_solve}'s \code{krylov\_step}---which is exactly where we go next.

\subsection{The richest contrast: \texttt{driven} and its per-\texorpdfstring{$\omega$}{omega} rebuild}
\label{sec:driven-contrast}

Set \code{driven} beside \code{electrostatic} and the family resemblance is total
except for one move. We render it more briefly, to surface the single difference:

\begin{lstlisting}[style=l4style]
driven :: DrivenConfig -> FrequencyResponse
driven cfg =
  let space  = nd_space cfg                                       -- H(curl) space
      fam    = { K  = fe_assemble space [ curl_curl (reluctivity cfg) ]  -- (1) assemble
               , C  = fe_assemble space [ damping cfg ]                  --     the basis
               , M  = fe_assemble space [ mass (permittivity cfg) ] }    --     ONCE (x3)
      omegas = sample_frequencies cfg                             -- the swept band
      es     = frequency_sweep fam omegas    -- (2) operator-VARYING map: A(w) rebuilt INSIDE
  in  sparameter_reduce (ports cfg) es       -- (3) port-projection -> S(omega)
\end{lstlisting}

\noindent The skeleton is identical: assemble, map, reduce. But three cells have
moved, exactly as \cref{ch:driven} promised. The \emph{assemble} now folds
\emph{three} term lists into a fixed \emph{basis} $\{\Kop,\Cop,\Mop\}$ rather than
one operator. The \emph{reduce} swaps \code{gram\_reduce} for
\code{sparameter\_reduce}, the port-projection (not a Gram). And the load-bearing
swap is the middle: \code{solve\_family} becomes \code{frequency\_sweep}. Both are
\texttt{map}s---independent solves, reorderable and parallel---but \code{solve\_family}
captures \emph{one fixed operator} outside the loop, while \code{frequency\_sweep}
rebuilds a \emph{fresh} operator $\mat{A}(\omega)=\Kop+i\omega\Cop-\omega^2\Mop$
\emph{inside} the loop, once per frequency, from the fixed basis. The rebuild lives in
\code{frequency\_sweep}'s body (\code{coordination}); the driver only composes it.
\Cref{fig:driven-dataflow} draws the per-$\omega$ rebuild loop.

\begin{figure}[t]
\centering
\begin{tikzpicture}[font=\footnotesize, >=Stealth, node distance=7mm and 10mm]
  \node[data, align=center, text width=15mm] (cfg) {\code{cfg}};
  \node[stage, right=of cfg, align=center, text width=22mm, fill=simBgGold, draw=simGold]
    (asm) {\code{fe\_assemble}$^{\times3}$\\[1pt]\scriptsize basis $\{\Kop,\Cop,\Mop\}$ once};
  \node[stage, right=of asm, align=center, text width=24mm, fill=simBgRust, draw=simRust]
    (sweep) {\code{frequency\_sweep}\\[1pt]\scriptsize per $\omega$:};
  \node[stage, right=of sweep, align=center, text width=20mm, fill=simBgGold, draw=simGold]
    (red) {\code{sparameter\_}\\\code{reduce}};
  \node[product, right=of red, text width=14mm, align=center] (S) {$S(\omega)$};
  \draw[flow] (cfg) -- (asm);
  \draw[flow] (asm) -- node[above,font=\scriptsize]{\code{fam}} (sweep);
  \draw[flow] (sweep) -- node[above,font=\scriptsize]{$[\vE_k]$} (red);
  \draw[flow] (red) -- (S);
  % the per-omega inner loop body
  \node[opbox, below=12mm of sweep, fill=white, align=center, text width=44mm] (inner)
    {\textbf{for each $\omega_k$:}\ form $\mat{A}(\omega_k)=\Kop{+}i\omega_k\Cop{-}\omega_k^2\Mop$\\
     $\to$ \code{ksp\_solve} (GMRES) $\to \vE_k$};
  \draw[refedge] (sweep.south) -- (inner.north);
  % the rebuild loop arrow
  \draw[backflow] (inner.east) .. controls +(1.4,0) and +(1.4,0) ..
    (inner.north east) node[midway, right=1mm, font=\scriptsize, text=simRust, align=center]
    {operator\\rebuilt\\\emph{inside}};
  \node[cornertag, above=0.5mm of asm]   {\code{data-algebra}};
  \node[cornertag, above=0.5mm of sweep] {\code{coordination}};
\end{tikzpicture}
\caption[The \texttt{driven} driver's per-\texorpdfstring{$\omega$}{omega} rebuild]{%
The \code{driven} driver's data flow, contrasting \code{electrostatic}
(\cref{fig:es-dataflow}). The \emph{basis} $\{\Kop,\Cop,\Mop\}$ is assembled once
(three \code{fe\_assemble} folds). But the system operator is \emph{not} fixed:
\code{frequency\_sweep} (the rust box) maps over the swept band, and its per-$\omega$
body (drawn below) \emph{rebuilds} $\mat{A}(\omega_k)$ from the basis and solves it
by GMRES---the operator is re-formed and re-captured \emph{inside} the loop (the rust
back-edge), the exact negation of the electrostatic hoist. The solved fields
$[\vE_k]$ feed \code{sparameter\_reduce}, the port-projection that yields $S(\omega)$.
One coordination shape (\code{solve\_family}) swapped for another
(\code{frequency\_sweep}); the rest of the skeleton is unchanged.}
\label{fig:driven-dataflow}
\end{figure}

\begin{notationbox}
\textbf{Magnetostatics: the same verb, a different weight.} The closest sibling of
all is \code{magnetostatic}. Its rendered body is \code{electrostatic} with the
$\Hone$ space changed to $\Hcurl$, the diffusion term changed to curl--curl, and
\emph{one closure changed}: the Gram weight becomes
\lstinline[style=l4style]|\i j -> 1 / (is!!i * is!!j)|, the current normalization
$1/(I_iI_j)$, where \code{is = currents cfg}. The solve corner is identical
(\code{solve\_family}, fixed-operator map); the reduction verb is identical
(\code{gram\_reduce}); only the weight closure differs. ``One verb, two products''
(\cref{ch:electrostatics}) is, in the rendering, one function \code{gram\_reduce}
called with two different weight lambdas.
\end{notationbox}

\section{The flagship: \texttt{krylov\_step} and CG as code}
\label{sec:cg-flagship}

We have seen a driver compose its solves; now we open one solve and watch the
abstract calculus produce a real algorithm. The \code{ksp\_solve} that
\code{solve\_family} maps---and that lives, in some form, inside \emph{every} driver
above---is at bottom a single step kernel, \code{krylov\_step}, folded by the
iteration combinator of \cref{ch:fold}. This is the showcase of the book: the place
where high-order combinators (\cref{ch:fold}), the state monad (\cref{ch:monad}),
and the strata of \cref{ch:strata} converge to synthesize the conjugate-gradient
method (\cref{ch:cg}) as runnable-looking code.

\subsection{Form A: the branch-in-body step}

The first rendering of \code{krylov\_step} is the direct one---a single step body
that does everything, with the first-iteration special case handled by a branch
inside it. It is the cleanest illustration of the monadic discipline of
\cref{ch:monad}: a sea of pure \code{let}s and exactly one \code{modify}.

\begin{lstlisting}[style=l4style]
-- Form A: branch-in-body. One step of CG, embedded in the Solve monad.
krylov_step :: OpParams -> Krylov -> (SimState -> Solve { krylov: Krylov, outputs: StepOutputs })
krylov_step op K = \s -> do
  -- ===== OUTSIDE the monad: pure let-bindings, the actual numerical work =====
  let Ap     = apply_linop op.A K.p     in   -- operator apply        (pure)
  let alpha  = K.beta / dot K.p Ap      in   -- step length           (pure dot)
  let x'     = axpy alpha K.p s.x       in   -- new iterate           (pure axpy)
  let r'     = axpy (negate alpha) Ap K.r in -- new residual          (pure axpy)
  let rdot'  = dot r' r'                in   -- residual energy       (pure dot)
  let beta'  = rdot' / K.beta           in   -- direction coefficient (pure)
  let p'     = axpy beta' K.p r'        in   -- new search direction   (pure axpy)
  let K'     = { ...K, r: r', p: p', beta: rdot' }  -- fresh workspace (plain value)
  let outputs = { residual_norm: sqrt (abs rdot') } -- demand-prunable readout
  -- ===== INSIDE the monad: the ONE state write =====
  modify (\s -> s { it: s.it + 1, x: x', converged: rdot' < op.tol })
  pure { krylov: K', outputs }
\end{lstlisting}

\noindent Count the monadic operations: there is exactly \emph{one}, the closing
\code{modify}. Everything above it---the operator apply \code{Ap}, the dot products
\code{alpha} and \code{rdot'}, the new iterate \code{x'}, residual \code{r'}, and
direction \code{p'}---is a pure \code{let}-binding, the algorithm's actual arithmetic
(\cref{ch:monad}). The three strata of \cref{ch:strata} are visible by eye:
\code{op} is read-only construction-stratum config (the operator, the tolerance);
\code{K} is the ephemeral per-restart workspace, threaded as a \emph{plain value} and
returned fresh; and the \code{SimState} \code{s} is the persistent run-time state,
touched only by the lone \code{modify}. The step \emph{returns} its readout in
\code{outputs} as a value, never logging it---so when no consumer reads the residual
history, the \code{residual\_norm} computation is pruned away (\cref{ch:fold}).

\subsection{Form B: first-iteration unrolling into bootstrap and steady}

Form~A carries a hidden cost. The conjugate-gradient recurrence builds each search
direction from the current residual \emph{and} a coefficient $\beta$ computed from
the \emph{previous} residual (\cref{ch:cg}). On the very first iteration there is no
previous residual, so the search direction is just $\vp\gets\vr$. Form~A would
handle that with an \lstinline[style=l4style]|if it == 0| branch inside the step---a
conditional that runs (and is mispredicted) on every iteration but is meaningful on
only the first.

The framework removes the branch by the \emph{first-iteration unrolling} rotation of
\cref{ch:fold}: split the step into a one-time \code{cg\_first\_step} that
manufactures the first $\beta$, followed by a branch-free \code{cg\_steady\_step}
that always has a real previous $\beta$ to work with, threaded by the combinator as a
\emph{closure argument} rather than a state field. This is
\code{iterate\_while\_with\_prev} (\cref{ch:fold}) in its defining use.
\Cref{fig:formAB} draws the rotation.

\begin{figure}[t]
\centering
\begin{tikzpicture}[font=\footnotesize, >=Stealth]
  % ===== LEFT: Form A, branch in body =====
  \begin{scope}
    \node[font=\small\bfseries, text=simRust] at (1.6,2.7) {Form A: branch in body};
    \node[stage, minimum width=34mm, minimum height=20mm, fill=simBgRust, draw=simRust,
          align=center] (stepA) at (1.6,1.0) {};
    \node[font=\scriptsize, anchor=north] at ($(stepA.north)+(0,-1mm)$) {\code{krylov\_step}};
    \node[diamond, draw=simRust, aspect=1.7, inner sep=1pt, font=\scriptsize,
          fill=white] (br) at (1.6,1.05) {\code{if it==0}};
    \node[font=\scriptsize\itshape, text=simRust, below=0.5mm of br] {tested every pass};
    \draw[backflow] (stepA.east) .. controls +(0.8,0) and +(0.8,0) .. (stepA.north east)
      node[midway, right=1mm, font=\scriptsize, text=simRust] {loop};
    \node[font=\scriptsize\itshape, text=simGray, align=center] at (1.6,-0.55)
      {one step, one\\extra slot \code{beta\_prev}};
  \end{scope}
  % rotation arrow
  \draw[depedge, line width=1pt] (3.7,1.0) -- node[above, font=\scriptsize, text=simBlue]
    {unroll} node[below, font=\scriptsize, text=simBlue] {iter.~1} (5.1,1.0);
  % ===== RIGHT: Form B, bootstrap + steady =====
  \begin{scope}[shift={(5.6,0)}]
    \node[font=\small\bfseries, text=simGreen] at (2.4,2.7) {Form B: bootstrap then steady};
    \node[opbox, minimum width=24mm, fill=simBgBlue, draw=simBlue, align=center]
      (boot) at (1.1,1.5) {\code{cg\_first\_step}\\\scriptsize runs once; makes $\beta_0$};
    \node[stage, minimum width=26mm, fill=simBgGreen, draw=simGreen, align=center]
      (steady) at (3.4,0.5) {\code{cg\_steady\_step}\\\scriptsize branch-free};
    \draw[flow] (boot.south) .. controls +(0,-0.5) and +(-1.2,0.4) ..
      node[below left=-1mm, font=\scriptsize, text=simBlue] {$\beta_0$} (steady.west);
    \draw[backflow] (steady.east) .. controls +(0.9,0) and +(0.9,0) .. (steady.north east)
      node[midway, right=1mm, font=\scriptsize, text=simGreen] {loop};
    \node[font=\scriptsize\itshape, text=simGreen, align=center] at (2.4,-0.7)
      {$\beta$ threaded as a CLOSURE arg\\(\code{iterate\_while\_with\_prev});\\carry one slot lighter};
  \end{scope}
\end{tikzpicture}
\caption[The Form A \texorpdfstring{$\to$}{->} Form B rotation]{The
first-iteration-unrolling rotation (\cref{ch:fold}) applied to \code{krylov\_step}.
\emph{Left, Form A}: a single step body carries an \lstinline[style=l4style]|if it==0|
branch (the $\beta_{k-1}$ has no value on iteration~0), tested on every pass though
meaningful only on the first, and a \code{beta\_prev} slot in the carry. \emph{Right,
Form B}: the loop is unrolled once into a \code{cg\_first\_step} (runs exactly once,
manufactures the first $\beta_0$) and a \emph{branch-free} \code{cg\_steady\_step};
the previous $\beta$ is threaded by \code{iterate\_while\_with\_prev} as a
\emph{closure argument}, so the steady carry is one slot lighter and the steady step
has no conditional. Same algorithm, rotated for a cleaner steady state.}
\label{fig:formAB}
\end{figure}

Here are the two unrolled steps and the driver that folds them, rendered in full.
This is the synthesized conjugate-gradient method:

\begin{lstlisting}[style=l4style]
-- the v0.5 CgState: ONE scalar lighter than Form A -- beta_prev is gone, threaded
-- instead as the PrevCarry closure parameter of iterate_while_with_prev.
type CgState = { x: Tensor[$S], r: Tensor[$S], p: Tensor[$S],
                 beta: Scalar, it: Int, converged: Bool }

-- the FIRST step: it == 0, so p <- r unconditionally (the Form-A branch, hoisted out)
cg_first_step :: LinOp[$S] -> Scalar -> CgState -> { state: CgState, residual_norm: Scalar }
cg_first_step opA eps s =
  let p'    = s.r                  in   -- it == 0  =>  p <- r   (no branch: this IS the first step)
  let Ap    = apply_linop opA p'   in
  let alpha = s.beta / dot Ap p'   in
  let x'    = axpy alpha p' s.x    in
  let r'    = axpy (negate alpha) Ap s.r in
  let beta' = dot r' r'            in
  let res'  = sqrt (abs beta')     in
  { state: { x: x', r: r', p: p', beta: beta', it: 1, converged: res' < eps },
    residual_norm: res' }

-- the STEADY step: branch-free; beta_prev is the closure carry (NOT a state field)
cg_steady_step :: LinOp[$S] -> Scalar -> Scalar -> CgState -> { state: CgState, residual_norm: Scalar }
cg_steady_step opA eps beta_prev s =
  let p'    = axpby 1.0 s.r (s.beta / beta_prev) s.p in   -- p <- r + (beta/beta_prev) p
  let Ap    = apply_linop opA p'   in
  let alpha = s.beta / dot Ap p'   in
  let x'    = axpy alpha p' s.x    in
  let r'    = axpy (negate alpha) Ap s.r in
  let beta' = dot r' r'            in
  let res'  = sqrt (abs beta')     in
  { state: { x: x', r: r', p: p', beta: beta', it: s.it + 1, converged: res' < eps },
    residual_norm: res' }

-- the DRIVER: first step, then fold cg_steady_step with iterate_while_with_prev,
-- threading the prior step's beta as the next step's beta_prev (closure, not a field).
cg_solve :: CgConfig -> LinOp[$S] -> Tensor[$S] -> Tensor[$S] -> Bool
         -> { final_state: CgState, residual_history: [Scalar] }
cg_solve config opA b x0 initial_guess =
  let { state: s0, initial_res } = cg_init opA b x0 initial_guess in
  let eps = max (config.rel_tol * initial_res) config.abs_tol in
  if sqrt (abs s0.beta) < eps                              -- already converged? do zero work
    then { final_state: { ...s0, converged: True }, residual_history: [] }
    else
      let { state: s1, residual_norm: res1 } = cg_first_step opA eps s0 in   -- the bootstrap
      if s1.converged || s1.it >= config.max_it
        then { final_state: s1, residual_history: [res1] }
        else
          let { final_state, trajectory } =
            iterate_while_with_prev
              (\_ -> pure { state: s1, prev: s0.beta })   -- bootstrap: seed (s1, beta_prev = s0.beta)
              s1                                          -- initial carry
              (\(s, beta_prev) ->                         -- steady body: (carry, prev)
                 let r = cg_steady_step opA eps beta_prev s in
                 pure { state: r.state, prev: s.beta, residual_norm: r.residual_norm })
              (\s -> s.it < config.max_it && not s.converged)   -- predicate: PURE on the carry
          in { final_state, residual_history: [res1] ++ trajectory.map (\t -> t.residual_norm) }
\end{lstlisting}

\noindent Read \code{cg\_solve} as the chapter's punchline. It is the
conjugate-gradient method of \cref{ch:cg}---a real, convergent linear solver---written
as a few lines over the combinator of \cref{ch:fold}. The \texttt{while} convention
of \cref{ch:fold} is visible in the very first \code{if}: an already-converged input
does \emph{zero} work. The bootstrap-then-steady split is
\code{iterate\_while\_with\_prev}. The predicate
\lstinline[style=l4style]|\s -> s.it < config.max_it && not s.converged| is pure on
the carry, reading the convergence flag and counter the step folded into the
state---never the residual the step just produced (\cref{ch:fold}'s predicate rule).
And the \code{residual\_history} is assembled from the trajectory the combinator
collected---demand-pruned away when nobody reads it. Everything those chapters built
abstractly is here, doing work.

\begin{keyidea}
The calculus produces real algorithm code. \code{cg\_solve} is the conjugate-gradient
method (\cref{ch:cg}) written as a short composition over
\code{iterate\_while\_with\_prev} (\cref{ch:fold}): a bootstrap step manufactures the
first $\beta$, a branch-free steady step folds the recurrence, the previous $\beta$
rides as a closure argument (\cref{ch:strata}), the predicate is pure on the carry,
the sole mutation is one \code{modify} of the counter (\cref{ch:monad}), and the
residual history is a demand-pruned trajectory. No part of it is pseudocode dressed
up---it is the abstract combinators of Part~III, instantiated, running a real solver.
\end{keyidea}

\section{Two worked examples}
\label{sec:worked}

The first worked example renders and walks the \code{electrostatic} driver against a
problem small enough to do by hand---from \code{IoData} projection to the capacitance
matrix---naming each library call as it fires. The second traces one CG step through
the synthesized \code{cg\_steady\_step}, tying every line to the conjugate-gradient
math of \cref{ch:cg}.

\begin{workedexample}{The \texttt{electrostatic} driver, run by hand}{es-driver}
We run the rendered driver of \S\ref{sec:es-rendered} on the two-conductor capacitor
of \cref{ch:electrostatics} (\cref{wex:es-twocond}), discretized to a single interior
degree of freedom, watching each library call do its job.

\medskip
\emph{(0) Project the config.} The parsed \code{IoData} has
\lstinline[style=l4style]|problemType cfg = ELECTROSTATIC|, so the lifecycle
(\S\ref{sec:lifecycle}) dispatched to \code{electrostatic} and handed it the config.
The driver reads its electrostatic view: \code{permittivity cfg} $=\varepsilon$, and
\lstinline[style=l4style]|terminal_sources cfg = [1, 2]| (two conductors). Line~(0)
builds \code{space = h1\_space cfg}, the scalar $\Hone$ space; on this coarse mesh it
has one free interior node.

\emph{(1) \code{fe\_assemble}, once.} Line~(1) folds the one-term list
\code{[\,diffusion}~$\varepsilon$\code{\,]} over the mesh. With $g=2\varepsilon A/L$
the half-gap conductance, the assembled-then-constrained stiffness is the scalar
$\Kop=2g$ (\cref{wex:es-twocond}). This is \code{k}, built a single time.

\emph{(2) The RHS family.} Line~(2)'s comprehension ranges over the two terminals.
\code{excitation cfg 1} encodes ``conductor~1 at one volt, conductor~2 grounded,''
giving $b_1=g$; \code{excitation cfg 2} gives $b_2=g$ by symmetry. So
\lstinline[style=l4style]|rhss = [g, g]|.

\emph{(3) \code{solve\_family k rhss}.} Line~(3) maps \code{ksp\_solve k} over the
family, reusing the \emph{one} operator \code{k} $=2g$ for both. Each solve is
$\Kop u=b$: $u_1=b_1/\Kop=g/(2g)=\tfrac12$ and $u_2=\tfrac12$. Reconstructing the full
nodal fields gives $V_1=(1,\tfrac12,0)$ and $V_2=(0,\tfrac12,1)$. So
\lstinline[style=l4style]|vs = [V1, V2]|---the solution family, produced by a map that
captured \code{k} once.

\emph{(4) \code{gram\_reduce k vs (\textbackslash i j -> 1)}.} Line~(4) pairs the
family against \code{k} with unit weight. The upper triangle is
$C_{11}=V_1^{\!\top}\Kop V_1=g/2$, $C_{22}=g/2$, and
$C_{12}=V_2^{\!\top}\Kop V_1=-g/2$; the verb mirrors $C_{21}=C_{12}$ (because $\Kop$
is symmetric). The returned product is
\[
  \mat{C}=\frac{g}{2}\begin{pmatrix}1&-1\\-1&1\end{pmatrix},
  \qquad -C_{12}=\frac{g}{2}=\frac{\varepsilon A}{L},
\]
the textbook parallel-plate capacitance. \tcblower
\emph{What the rendering added.} The arithmetic is identical to
\cref{wex:es-twocond}---but now we watched it run as \emph{code}. The same \code{k}
appeared in line~(3) (the solve) and line~(4) (the reduce): one value, two
consumers, the back-edge of \cref{fig:es-dataflow} made literal. The ``two solves''
were the two elements of \code{rhss}, not a hand-written loop. And the whole thing had
type \code{ElectrostaticConfig -> CapacitanceMatrix}: \code{IoData} in, $\mat{C}$ out,
nothing mutable crossing the boundary. On a real mesh only line~(3) changes
character---the two scalar divisions become two million-dimensional CG solves, each
the \code{cg\_solve} of \S\ref{sec:cg-flagship}---but the composition is exactly the
five lines we ran.
\end{workedexample}

\begin{workedexample}{One CG step through \texttt{cg\_steady\_step}}{cg-step}
We trace a single iteration of the synthesized \code{cg\_steady\_step}
(\S\ref{sec:cg-flagship}) and tie each line to the conjugate-gradient math of
\cref{ch:cg}. Take the tiny SPD system
\[
  \mat{A}=\begin{pmatrix}4&1\\1&3\end{pmatrix},\quad \vb=\begin{pmatrix}1\\2\end{pmatrix},
\]
and suppose the loop has reached a steady state with carry
\lstinline[style=l4style]|s = { x, r, p, beta, it, ... }| where
$\vr=(\tfrac{1}{2},\tfrac{1}{2})^{\!\top}$, $\vp=(1,0)^{\!\top}$, the carry's
$\code{beta}=\inner{\vr}{\vr}=\tfrac12$, and the closure carry
$\code{beta\_prev}=2$ (the previous step's residual energy). We run
\code{cg\_steady\_step opA eps beta\_prev s} line by line.

\medskip
\emph{Line 1 --- the direction update \code{p' = axpby 1.0 s.r (s.beta/beta\_prev) s.p}.}
The coefficient is $\beta/\beta_{\text{prev}}=\tfrac12/2=\tfrac14$. This is the CG
$\beta_k=\inner{\vr_k}{\vr_k}/\inner{\vr_{k-1}}{\vr_{k-1}}$ of \cref{ch:cg}, and the
threaded \code{beta\_prev} is exactly $\inner{\vr_{k-1}}{\vr_{k-1}}$---the
\emph{previous} residual energy, supplied as a closure argument, not read from a state
field. The new direction is
$\vp'=\vr+\tfrac14\vp=(\tfrac12+\tfrac14,\ \tfrac12)=(\tfrac34,\tfrac12)^{\!\top}$.

\emph{Line 2 --- the apply \code{Ap = apply\_linop opA p'}.} One operator application,
$\mat{A}\vp'=(4\cdot\tfrac34+\tfrac12,\ \tfrac34+3\cdot\tfrac12)=(\tfrac72,\tfrac94)^{\!\top}$.
This is the single matrix--vector product that is the cost of a CG step (\cref{ch:cg});
in the driver it is one \code{krylov\_step}'s operator apply, the matrix-free leaf of
\cref{ch:fulltrace}.

\emph{Line 3 --- the step length \code{alpha = s.beta / dot Ap p'}.} The denominator
is $\inner{\mat{A}\vp'}{\vp'}=\tfrac72\cdot\tfrac34+\tfrac94\cdot\tfrac12=\tfrac{21}{8}+\tfrac98=\tfrac{30}{8}=\tfrac{15}{4}$,
so $\alpha=\beta/\inner{\mat{A}\vp'}{\vp'}=\tfrac12/\tfrac{15}{4}=\tfrac{2}{15}$.
This is the CG line-search $\alpha_k=\inner{\vr_k}{\vr_k}/\inner{\vp_k}{\mat{A}\vp_k}$
(\cref{ch:cg}), the step that minimizes the energy norm along $\vp'$.

\emph{Lines 4--5 --- the updates \code{x' = axpy alpha p' s.x},
\code{r' = axpy (-alpha) Ap s.r}.} The iterate moves $\vx'=\vx+\alpha\vp'$ and the
residual updates $\vr'=\vr-\alpha\mat{A}\vp'$---the two \code{axpy} fused-multiply-adds
that advance CG (\cref{ch:cg}). Numerically
$\vr'=(\tfrac12,\tfrac12)-\tfrac{2}{15}(\tfrac72,\tfrac94)=(\tfrac12-\tfrac{7}{15},\ \tfrac12-\tfrac{3}{10})=(\tfrac{1}{30},\tfrac{1}{5})^{\!\top}$.

\emph{Lines 6--7 --- the new residual energy and convergence.}
$\code{beta'}=\inner{\vr'}{\vr'}=(\tfrac{1}{30})^2+(\tfrac15)^2\approx0.0411$, and
$\code{res'}=\sqrt{\code{beta'}}\approx0.203$---a smaller residual than the
$\sqrt{\code{beta}}=\sqrt{\tfrac12}\approx0.707$ we started with. The step writes
\code{converged: res' < eps} into the state.

\tcblower
\emph{Reading the trace as the calculus.} Every CG quantity of \cref{ch:cg} appears
exactly once and in the expected place: $\beta_k$ on line~1, the apply on line~2,
$\alpha_k$ on line~3, the iterate and residual updates on lines~4--5. The step is
\emph{branch-free}---there is no \lstinline[style=l4style]|if it==0|---because the
unrolling of \S\ref{sec:cg-flagship} moved that case into \code{cg\_first\_step}; and
the previous residual energy reached this step as the closure argument
\code{beta\_prev}, threaded by \code{iterate\_while\_with\_prev} (\cref{ch:fold}), not
stored in \code{CgState}. The one effect on the persistent \code{SimState}---the
counter bump and convergence flag---is the lone \code{modify} of Form~A
(\cref{ch:monad}); the iterate, residual, and direction are pure \code{let}s. This is
the conjugate-gradient method of \cref{ch:cg}, running, with each line a calculus
construct we built earlier in the book.
\end{workedexample}

\section{Writing a new analysis, and the spine that runs them}
\label{sec:lifecycle}

We close by stepping back from the six existing drivers to the shape they share, and
asking what it would take to write a seventh.

\subsection{How a new analysis is written}

Because a driver is nothing but a composition, adding an analysis is a matter of
\emph{three choices}, not a new program (\cref{fig:new-analysis}):

\begin{enumerate}[nosep]
  \item \textbf{Pick the corner.} Decide how the solves coordinate: a fixed-operator
  map (\code{solve\_family}), an operator-varying map (\code{frequency\_sweep}), a
  state-threaded fold (\code{fold\_solve}), or a single opaque call (\code{eigsolve}).
  This is the choice of coordination shape from \code{coordination}.
  \item \textbf{Pick the reduction.} Decide what physical product the solved family
  collapses to: a symmetric Gram (\code{gram\_reduce}), a port-projection
  (\code{sparameter\_reduce}), a per-mode table (\code{eigenfreq\_qfactor\_reduce}),
  an energy table (\code{domain\_energy\_reduce}). This is a verb from
  \code{data-algebra}.
  \item \textbf{Compose.} Wire \code{fe\_assemble} (the right term list) into the
  chosen corner into the chosen reduction. That composition \emph{is} the driver.
\end{enumerate}

\noindent No new combinator, no new record, no new loop. The expressive economy is
the whole argument of the book cashed out: the simulator is a small library, and an
analysis is a short composition over it. The four existing reduction verbs and four
coordination shapes already span the six analyses Palace ships; a genuinely new
analysis would reuse them and differ only in its term list and its choice of corner
and reduction.

\begin{figure}[t]
\centering
\begin{tikzpicture}[font=\footnotesize, >=Stealth, node distance=5mm]
  \node[data, align=center, text width=20mm] (cfg) at (0,0) {term list +\\config view};
  \node[stage, fill=simBgGold, draw=simGold, right=10mm of cfg, text width=18mm, align=center]
    (asm) {\code{fe\_assemble}\\\scriptsize \code{data-algebra}};
  % the corner choice
  \node[stage, fill=simBgRust, draw=simRust, right=9mm of asm, text width=24mm, align=center]
    (corner) {\textbf{pick the corner}\\\scriptsize map / map / fold / call\\\code{coordination}};
  \node[stage, fill=simBgBlue, draw=simBlue, right=9mm of corner, text width=22mm, align=center]
    (red) {\textbf{pick the reduce}\\\scriptsize Gram / proj / table\\\code{data-algebra}};
  \node[product, right=7mm of red, text width=14mm, align=center] (prod) {new\\product};
  \draw[flow] (cfg)--(asm); \draw[flow] (asm)--(corner);
  \draw[flow] (corner)--(red); \draw[flow] (red)--(prod);
  \node[cornertag, below=1mm of corner] {choice 1};
  \node[cornertag, below=1mm of red] {choice 2};
  \node[font=\scriptsize\itshape, text=simGray] at (6.3,-1.5)
    {choice 3: compose the three into one driver def};
\end{tikzpicture}
\caption[Writing a new analysis]{Writing a new analysis is three choices, not a new
program. Start from a weak-form term list and a config view; assemble with
\code{fe\_assemble} (\code{data-algebra}); \textbf{pick the corner}---one of the four
coordination shapes (\code{coordination}); \textbf{pick the reduction}---one of the
verbs (\code{data-algebra}); and \textbf{compose} the three into a driver definition.
No new combinator or record is written; the new analysis reuses the same library and
differs only in its term list and its two choices. This is why the six drivers are
each a handful of lines.}
\label{fig:new-analysis}
\end{figure}

\subsection{The spine root: \texttt{lifecycle} dispatches over the six}

One definition sits above all six drivers and runs them: the \code{lifecycle} root
(\cref{ch:fulltrace}). It is the topologically last definition in the library---it
\emph{composes the driver defs themselves}---and it is short, because dispatch is one
field read:

\begin{lstlisting}[style=l4style]
lifecycle :: IoData -> Product
lifecycle cfg =
  let mesh0 = build_mesh cfg                          -- shared front end: load + partition mesh
      drv   = dispatch (problemType cfg)              -- the specialization seam: pick ONE driver
      step  = \m -> estimate_mark_refine cfg (drv cfg m)  -- run the driver, then mark/refine
  in  fold_solve_amr step mesh0 (refinement cfg)      -- adaptive estimate-mark-refine OUTER fold
  where
    dispatch ELECTROSTATIC = electrostatic            -- the six driver defs, by problem type
    dispatch MAGNETOSTATIC = magnetostatic
    dispatch DRIVEN        = driven
    dispatch TRANSIENT     = transient
    dispatch EIGENMODE     = eigenmode
    dispatch BOUNDARYMODE  = boundary_mode
\end{lstlisting}

\noindent Three things to read. First, \code{build\_mesh} and the parse that produced
\code{cfg} are the \emph{shared front end} (\cref{ch:fulltrace})---driver-agnostic,
run identically whichever analysis follows. Second, \code{dispatch (problemType cfg)}
is the \emph{specialization seam}: one field of the one \code{IoData} selects one of
the six driver defs above. Everything before the seam is shared; the seam picks the
one stage that differs. Third, the whole run is wrapped in \code{fold\_solve}---the
adaptive estimate--mark--refine loop is itself a fold (\cref{ch:fold}), whose carry is
the mesh, degenerating to a single solve when refinement is off. The spine root is the
\emph{fourth} use of the same coordination vocabulary: a fold wrapped around a
dispatch wrapped around the six compositions.

\begin{keyidea}
The whole simulator is one library and a one-field dispatch. \code{lifecycle} parses
\code{IoData}, builds the mesh (the shared front end), reads \code{problemType} to
select one of the six driver compositions (the specialization seam), and threads the
adaptive loop through \code{fold\_solve}. Above the six few-line drivers sits one
few-line root that runs any of them. From parse to product, the entire machine is the
five-module library of \cref{ch:fivelibs} composed---no glue beyond the compositions
themselves.
\end{keyidea}

\subsection{Where this leaves us}

We set out, in \cref{ch:fulltrace}, to \emph{draw} the simulator as one composition;
in this chapter we \emph{wrote} it. The six drivers are six short compositions over
four libraries; the conjugate-gradient solver inside every one of them is a few lines
over the combinator of \cref{ch:fold}; and a seventh analysis would be three choices
and a composition. The synthesized library is not a metaphor for the simulator---it is
the simulator, with its compositions made into code. What remains is to step back one
final time and see the whole machine---library, drivers, and lifecycle---as a single
object, which is the closing chapter, \cref{ch:wholemachine}.

\summarysection
The \code{drivers} library is the top module of the synthesized library
(\cref{ch:fivelibs}): it holds the six entry-point analyses, each a short composition
over the four calculus libraries (\code{types}, \code{iteration}, \code{data-algebra},
\code{coordination}) and adding no new vocabulary. The six are one-liners in a single
algebra---\code{electrostatic} $=$ \code{gram\_reduce}$(w{=}1)\circ$%
\code{solve\_family}$\circ$\code{fe\_assemble}; \code{magnetostatic} the same with the
weight $1/(I_iI_j)$; \code{driven} $=$ \code{sparameter\_reduce}$\circ$%
\code{frequency\_sweep}$\circ$\code{fe\_assemble}$^{\times3}$; \code{transient} $=$
\code{fold\_solve}$\circ$\code{fe\_assemble}$^{\times3}$; \code{eigenmode} and
\code{boundary\_mode} $=$ \code{map readout}$\circ$\code{eigsolve}$\circ$%
\code{eig\_pencil}$\circ\cdots$. Beneath the four solve corners of \cref{ch:situating}
lie just three coordination shapes---\texttt{map} (independent solves), \texttt{fold}
(sequential chain), and \texttt{single call} (opaque eigensolve)---and every driver is
one of the three wrapped around \texttt{assemble} and capped by \texttt{reduce}. The
per-driver config types are projection-views of the one \code{IoData}, not new shapes.
We rendered \code{electrostatic} in full and walked its five lines, naming the library
each call comes from---\code{fe\_assemble} assembles $\Kop$ once (\code{data-algebra});
a comprehension builds the per-terminal RHS family; \code{solve\_family} maps the
fixed-operator solve (\code{coordination}); \code{gram\_reduce} pairs the family
against the same $\Kop$ (\code{data-algebra})---and contrasted \code{driven}, whose only
real difference is rebuilding $\mat{A}(\omega)$ \emph{inside} \code{frequency\_sweep}.
The flagship is \code{krylov\_step}/CG: Form~A (branch-in-body, one \code{modify} amid
pure \code{let}s) rotated by first-iteration unrolling into Form~B
(\code{cg\_first\_step} $+$ branch-free \code{cg\_steady\_step}), folded by
\code{iterate\_while\_with\_prev} in \code{cg\_solve}---the conjugate-gradient method of
\cref{ch:cg} as runnable code, with $\beta$ threaded as a closure argument
(\cref{ch:strata}), the predicate pure on the carry (\cref{ch:fold}), and the residual
history demand-pruned. A new analysis is three choices---pick the corner, pick the
reduction, compose---and the \code{lifecycle} root parses \code{IoData}, builds the
mesh, dispatches on one field to one of the six drivers, and wraps the adaptive loop in
\code{fold\_solve}. Every driver is a few lines of map/fold/single-call over
\texttt{assemble}$\to$\texttt{reduce}, and the abstract calculus of Part~III produces
real algorithm code.

\furtherreading
The conjugate-gradient method rendered here as \code{cg\_solve}, and the broader family
of Krylov solvers the other drivers invoke, are developed in full by
Saad~\citep{saad2003}, whose treatment of CG, its short recurrence, and the role of the
$\beta$ coefficient is the standard reference behind \cref{ch:cg} and this chapter's
flagship. The functional-programming machinery that makes the synthesized rendering
possible---the state monad threading \code{SimState}, the \code{do}-notation, the
higher-order combinators that \code{iterate\_while\_with\_prev} specializes---is
documented by Peyton Jones~\citep{peytonjones2003}, the source behind \cref{ch:monad}
and \cref{ch:fold}. The simulator whose six drivers we have rendered is Palace itself,
documented at the project~\citep{palace2023}; readers comparing the synthesized
compositions to the original C++ will find the driver dispatch and the per-analysis
solvers there, in the imperative form this book has spent six parts rotating into the
pure compositions of this chapter. The closing chapter, \cref{ch:wholemachine},
assembles library, drivers, and lifecycle into one final view of the whole machine.

\exercisessection
\begin{enumerate}[nosep]
  \item \textbf{Read the one-liners.} For each of the six driver one-liners (the list
  at the opening of the chapter), name (a) how many operators \code{fe\_assemble}
  builds, (b) which coordination shape runs the solve (\code{solve\_family} /
  \code{frequency\_sweep} / \code{fold\_solve} / \code{eigsolve}), and (c) which
  reduction caps it. Which two drivers share the \emph{same} coordination shape but
  differ only in a weight?
  \item \textbf{Walk the rendering.} Take the rendered \code{electrostatic} driver of
  \S\ref{sec:es-rendered} and, for each of its five body lines, name the library the
  call comes from and the one fact the line establishes. Which value is \code{let}-bound
  on one line and consumed on \emph{two} later lines, and why does that double use
  appear as a back-edge in \cref{fig:es-dataflow}?
  \item \textbf{The one moved call.} Comparing \code{electrostatic} and \code{driven},
  identify the single coordination call that swaps (\code{solve\_family} for
  \code{frequency\_sweep}) and explain, in terms of \emph{where} the operator capture
  sits relative to the loop, why one rebuilds the operator inside the map and the other
  does not. Cross-check against \cref{ch:driven}.
  \item \textbf{Magnetostatics as one closure.} Magnetostatics calls the \emph{same}
  \code{gram\_reduce} as electrostatics. Write the weight closure each passes, and
  explain why the inductance matrix and the capacitance matrix come from one verb with
  two lambdas rather than two reduction routines. (See the notation box in
  \S\ref{sec:driven-contrast}.)
  \item \textbf{Count the monad.} In Form~A of \code{krylov\_step}
  (\S\ref{sec:cg-flagship}), classify every line as ``in the monad'' or ``outside the
  monad'' by the rule of thumb of \cref{ch:monad}. How many are in the monad? Name the
  three strata (\cref{ch:strata}) the step touches and which variable carries each.
  \item \textbf{Why unroll.} Explain what the first-iteration unrolling
  (\cref{fig:formAB}) buys: what branch does \code{cg\_steady\_step} avoid that Form~A
  carries, where did that case go, and how is the previous $\beta$ supplied to the
  steady step if it is no longer a field of \code{CgState}? Which combinator threads it?
  \item \textbf{Trace a CG step.} Repeat \cref{wex:cg-step} for the \emph{next} step:
  use the carry it produced (the new $\vr'$, $\vp'$, and $\beta'=\inner{\vr'}{\vr'}$),
  take \code{beta\_prev} to be the previous step's $\beta=\tfrac12$, and run
  \code{cg\_steady\_step} line by line to compute the new direction, step length,
  iterate, and residual. Confirm the residual norm continues to fall, and name the CG
  formula each line implements.
  \item \textbf{Write the seventh analysis.} Suppose you wanted a driver that computes
  the \emph{energy} stored by each of several independent excitations (no cross terms).
  Using the three-choice recipe of \S\ref{sec:lifecycle}: which coordination shape would
  you pick, which reduction verb, and what term list would \code{fe\_assemble} take?
  Sketch the five-line driver by analogy with \code{electrostatic}, and say which one
  line differs from it.
\end{enumerate}
